\chapter{AI 使用说明与复现实验}

\section{AI 在开发过程中的位置}

这轮工作使用 AI 辅助阅读跨组件调用链、归纳日志、讨论候选方案、补充定向测试和整理文档。由于性能问题很容易被局部数据误导，AI 输出不能直接作为实现结论。热点是否真实、锁能否缩小、缓存对象是否稳定，以及一项改动是否进入主线，都由源码审查和实测决定。

在 syscall/VFS 优化中，AI 帮助把 syscall 入口、用户拷贝、fd registry、VFS session 和具体文件系统的重复访问连成一条路径，并提示检查阻塞点与文件偏移的线性化边界。在 MM 和零页池优化中，AI 用于枚举 lazy VMA、部分失败、远端 TLB、OOM drain、锁顺序和 PPN 所有权等边界。实验参数、回退门槛和最终合入仍以维护者判断为准。

AI 也会提出无效方案。page cache reverse index、child index、TLSF 小对象池和扩大预读等方向都曾在局部上看起来合理，但完整轮没有收益或发生正确性问题。保留这些反例，是为了说明我们没有把生成的建议当成既定路线，也没有只挑选对结论有利的数据。

\section{复现环境与证据分级}

建议使用与记录相同的 QEMU 9.2.1、RISC-V \code{virt}、8 vCPU、16 GiB guest 内存、final 内核、同一赛题镜像和 \code{-snapshot}。A/B 两侧只改变待测提交或 feature，不能在中途更换镜像、QEMU 参数、宿主 CPU 集合或 workload。对于跨架构公共代码，还应在同一提交上完成 RISC-V 和 LoongArch 检查。

实验记录按证据强度分为三类：

\begin{enumerate}
  \item \textbf{完整 A/B}：具有相同条件、完整成功标记和相邻对照，可用于判断墙钟收益；
  \item \textbf{固定窗口诊断}：pc-hot、wait-hot 或 syscall profile 到达固定时限，用于解释热点与机制；
  \item \textbf{机制统计}：命中率、变化摘要、锁等待或 refill 来源，用于确认代码按设计工作。
\end{enumerate}

后两类不能单独写成“性能提升”。如果固定窗口中候选执行了更多指令，可能只是推进了更多 workload；如果缓存命中率很高，也仍可能因锁竞争和淘汰开销而变慢。

\clearpage
\section{pc-hot 与 syscall profile 的复现方法}

pc-hot 需要与内核匹配的 ELF 符号。下面给出最小流程，QEMU 后面的磁盘和网络参数应与正式 runner 保持一致：

\begin{verbatim}
cd os
./scripts/pc-hot/pc-hot-rv.sh build
./scripts/pc-hot/pc-hot-rv.sh run /tmp/pcs-rv.txt -- \
  timeout 300 qemu-system-riscv64 -machine virt \
  -kernel ./kernel-rv-final -m 16G -smp 8 -nographic ...
./scripts/pc-hot/pc-hot-rv.sh analyze \
  /tmp/pcs-rv.txt ./kernel-rv-final 50
\end{verbatim}

分析候选和对照时，输出文件必须分开，ELF 也不能混用。除了 top symbols，还应记录总指令、每个 vCPU 分布和 workload 到达位置。需要区分计算与等待时，可在相同参数下换用 \code{wait-hot-rv.sh}。

syscall profile 的 RISC-V full-system 入口如下：

\begin{verbatim}
./scripts/syscall-profile/syscall-profile-rv.sh build
./scripts/syscall-profile/syscall-profile-rv.sh run \
  /tmp/syscalls.txt backend=auto paths=1 \
  max_path=256 top_paths=200 -- \
  timeout 300 qemu-system-riscv64 -machine virt \
  -kernel ./kernel-rv-final -m 16G -smp 8 -nographic ...
./scripts/syscall-profile/analyze.py \
  /tmp/syscalls.txt --top 30 > /tmp/syscalls.md
\end{verbatim}

full-system 下应确认 \code{backend=auto} 实际选择 \code{ecall}。强制使用 QEMU native syscall callback 在 WaterOS full-system 实测为零事件，因为 syscall 由 guest 内核处理；该后端主要用于 linux-user 对照。

\section{完整 BuildStorm A/B}

完整复现可按以下顺序进行：

\begin{enumerate}
  \item 为 baseline 和 candidate 分别记录 \code{git rev-parse HEAD}，构建 final kernel 并保存 SHA-256；
  \item 执行 \code{make rv\_check} 与 \code{make la\_check}，再运行改动涉及的功能测例；
  \item 固定宿主 CPU 亲和性、QEMU 9.2.1、镜像哈希、SMP 与内存，确认磁盘使用 \code{-snapshot}；
  \item 交错运行完整 A/B，至少取得两轮候选和相邻对照；
  \item 检查 toolchain、minibuild 和 multi compile 标记，不只提取最后一个数字；
  \item 计算各自中位数与相对变化，并把 timeout、panic、SIGSEGV、OOM 或缺项记为失败，而不是删除样本；
  \item 对超过门槛的候选再运行 pc-hot 或机制统计，确认没有把开销转移到其他子系统。
\end{enumerate}

syscall-chain 任务将大改动拆成可回滚的小提交：测量 runner、mount namespace、dentry、FD slot、I/O lease、lookup token、exec 稳定句柄、lazy ELF、mmap/munmap 变化摘要、只读路径和 futex 分片分别验收。复现时可沿这些提交定位行为变化，但不能把只有静态检查、尚无 QEMU A/B 的任务写成已确认性能收益。

\section{文档和实现如何对应}

性能任务索引记录了早期 benchmark 缺口、风险级别和 FS/MM/allocator 的候选；独立实验文档保存工具设计、固定条件、命中率和回退原因；syscall-chain 简报保存每个接口改动的语义边界。\code{perf/*} 分支则保留代码和失败实验。本文把这些材料按因果链重新组织，数字仍应回到对应实验记录核对。

若后续修改正文中的结果，至少同时更新：被测提交与参数、结果表中的证据类型、失败样本、最终内核上的重新验证状态。这样可以避免文档随着主线演进而留下无法复现的“历史最好值”。
